\documentclass{situationist-preprint}

\title{Surprise and Inevitability}
\msid{MTH-13-a}
\msfamily{MTH-13}
\msdiscipline{Philosophy of Mathematics}
\mstype{Principal manuscript}
\mskeywords{mathematical beauty; surprise; inevitability; response-dependence; Hardy; explanation; aesthetics of proof}

\begin{document}
\maketitle

\begin{abstract}
Mathematicians call theorems and proofs beautiful and mean it as praise of a distinctive kind, not reducible to truth, depth or usefulness. We argue that mathematical beauty is the combination of surprise and inevitability: a proof is beautiful when it leads a qualified observer to a conclusion she could not have anticipated from what she understood beforehand, by a route that, once grasped, shows that the conclusion could not have been otherwise. Theorems are beautiful derivatively, through the proofs that make them inevitable. The account is response-dependent, as beauty in general is: the surprise and the inevitability are relative to an observer who understands the statement and its background, and the beauty consists in the disposition to produce, in such observers, the pleasure of grasping both at once. We show that neither component suffices alone, that the two are compatible because they are indexed to different stages of understanding, and that the account explains why non-explanatory proofs, among them most computer-assisted ones, are seldom called beautiful. Simplicity, enlightenment and unity in variety are partial versions of it.
\end{abstract}

\section{The question}

\citet{Hardy1940} held that the best mathematics is beautiful as well as serious, and he analysed the beauty of a proof into three qualities: unexpectedness, inevitability and economy. \citet{Poincare1908} described an aesthetic sensibility, attuned to harmony and elegance, which selects among the combinations of ideas that the mind produces and lets the useful ones into consciousness. \citet{Hutcheson1725}, two centuries earlier, had counted theorems among the objects of the sense of beauty, and explained their beauty, as he explained beauty in general, by uniformity amidst variety: a single theorem comprehends an infinity of particular cases. \citet{Rota1997} argued that ascriptions of beauty in mathematics are an oblique way of acknowledging enlightenment, the sense of understanding why a result holds, and that the aesthetic vocabulary is a misdescription of a cognitive achievement.

Recent work has been both empirical and historical. \citet{InglisAberdein2015} surveyed mathematicians' aesthetic judgments of a proof and found that they vary along several independent dimensions, so that beauty is not simply a matter of simplicity. \citet{ZekiEtAl2014} found that mathematicians' experience of formulae they rate as beautiful engages the same region of the brain as the experience of beauty in other domains. \citet{Breitenbach2015} reconstructed a Kantian account on which the beauty of a proof lies in the free play between the understanding and the imagination it occasions, and \citet{Montano2014} developed a naturalistic aesthetic theory of mathematics.

We defend Hardy's first two qualities as the whole of the account. Mathematical beauty is the combination of surprise and inevitability. Economy, we shall argue, matters only as a condition of inevitability's being grasped. The account is response-dependent in the sense of AES-02-a, on which beauty in general is a disposition of objects to give a distinctive pleasure to qualified observers in suitable conditions. What is specific to mathematics is the kind of pleasure and the kind of observer.

\section{Framework}

\begin{definition}[Qualified observer]\label{d:obs}
A \emph{qualified observer} of a mathematical proof is one who understands the statement of its conclusion and the concepts it involves, knows the background against which the result is posed, including what was antecedently known and expected, and follows the proof step by step. \emph{Suitable conditions} are those in which she attends to the proof as a whole, without haste and without prior familiarity with its route.
\end{definition}

The definition specializes the qualified observer of AES-02-a to mathematics: delicacy of discrimination becomes understanding, and acquaintance with comparable objects becomes knowledge of the background.

\begin{definition}[Surprise]\label{d:surp}
A proof's conclusion, or the connection by which the proof reaches it, is \emph{surprising} for a qualified observer if she could not have anticipated it from her understanding of the statement and its background before following the proof.
\end{definition}

\begin{definition}[Inevitability]\label{d:inev}
A proof makes its conclusion \emph{inevitable} for a qualified observer if, once she has followed it, she sees that the conclusion could not have been otherwise, because it follows from a feature of the objects concerned that she now holds in view as a whole.
\end{definition}

\begin{principle}[Response-dependence]\label{p:rd}
A mathematical object of appreciation is beautiful if and only if it is disposed to give a distinctive pleasure to qualified observers in suitable conditions, where the observers and conditions are specified without using the concept of beauty.
\end{principle}

\begin{principle}[Discrimination]\label{p:discr}
An account of mathematical beauty must distinguish it from truth, depth, importance and difficulty, each of which a result can have without being beautiful.
\end{principle}

Principle~\ref{p:rd} is the thesis of AES-02 applied to mathematics. Principle~\ref{p:discr} records a datum of practice: the classification of the finite simple groups is deep, important and difficult, and is not usually called beautiful. Inevitability in Definition~\ref{d:inev} is not mere validity. Every valid proof makes its conclusion follow; only some let the observer see that it must, because they derive it from a feature she can hold in view rather than through a sequence of steps she can only check one at a time.

\section{The account}

\begin{theorem}\label{t:main}
A proof is mathematically beautiful if and only if it is disposed to give qualified observers in suitable conditions the pleasure of grasping its conclusion as both surprising and inevitable. A theorem is beautiful if and only if it has a beautiful proof.
\end{theorem}

The argument has three parts: each component is necessary, the two together are sufficient, and they are compatible.

\begin{lemma}\label{l:surp}
Surprise without inevitability is not beauty.
\end{lemma}

\begin{proof}
The number $e^{\pi\sqrt{163}}$ differs from an integer by less than $10^{-12}$. The fact is surprising for any qualified observer; established by computation, it is not beautiful but curious, a numerical coincidence. It becomes an object of mathematical admiration when it is derived from the theory of complex multiplication, which shows that the near-integrality follows from the fact that the imaginary quadratic field of discriminant $-163$ has class number one. That derivation supplies inevitability, and with it the response that Principle~\ref{p:rd} identifies with beauty. Since the surprise was present before and the beauty was not, surprise is not sufficient.
\end{proof}

\begin{lemma}\label{l:inev}
Inevitability without surprise is not beauty.
\end{lemma}

\begin{proof}
A routine verification, such as the computation of a derivative by the rules of differentiation or the verification of an identity by expanding both sides, makes its conclusion inevitable, in that the observer sees that it follows from the definitions she holds in view. But she expected it, and could have anticipated both the conclusion and the route. Such verifications are correct and may be elegant in execution, but they are not called beautiful. So inevitability is not sufficient.
\end{proof}

\begin{proposition}\label{p:compat}
Surprise and inevitability are compatible.
\end{proposition}

\begin{proof}
They are indexed to different stages of the observer's understanding. Surprise is relative to her understanding before following the proof (Definition~\ref{d:surp}), inevitability to her understanding after (Definition~\ref{d:inev}). A conclusion that could not have been anticipated from the statement can be seen to follow necessarily from a feature the proof brings into view, and this is exactly what happens when a proof reveals a connection the observer did not know.
\end{proof}

\begin{proof}[Proof of Theorem~\ref{t:main}]
By Lemmas~\ref{l:surp} and~\ref{l:inev}, each component is necessary, since its absence leaves a case that is not beautiful. For sufficiency, consider the proofs that mathematicians most often cite as beautiful. Euclid's proof that there are infinitely many primes reaches its conclusion by an unanticipated construction, the product of the given primes plus one, and once the construction is seen the conclusion cannot fail, because the new number must have a prime factor not on the list. Cantor's diagonal argument surprises, since nothing in the statement suggests that infinite sets differ in size, and makes the conclusion inevitable, since the diagonal sequence differs from every listed one by construction. Euler's evaluation of $\sum 1/n^2$ as $\pi^2/6$ joins a sum over the integers to the geometry of the circle, which no observer expects, and the derivation from the factorization of the sine function shows why $\pi$ must appear. In each case, the response of qualified observers combines the two, and the cases are those in which the judgment of beauty is most secure. By Principle~\ref{p:rd}, the disposition to produce that response is the beauty. That a theorem is beautiful when it has a beautiful proof follows, since what is admired in a theorem is the inevitability of a surprising result, which only a proof can supply.
\end{proof}

The account satisfies Principle~\ref{p:discr}. Truth is shared by every theorem. Depth and importance concern a result's connections with the rest of mathematics, and difficulty the effort it costs; a deep, important and difficult result proved by exhausting thousands of cases is surprising in its conclusion but not inevitable in its proof.

\begin{corollary}\label{c:expl}
A proof that does not explain why its conclusion holds is not beautiful with respect to the part that does not explain. In particular, computer-assisted proofs are generally not beautiful.
\end{corollary}

\begin{proof}
MTH-08 argues that a proof explains when it derives its conclusion uniformly from a feature salient in the theorem, and that computer-assisted proofs generally do not, because they exhaust cases separately. A proof that exhausts cases separately does not let the observer see that the conclusion follows from a feature she holds in view as a whole, so by Definition~\ref{d:inev} it does not make the conclusion inevitable, and by Theorem~\ref{t:main} it is not beautiful.
\end{proof}

The corollary agrees with the reception of the four colour theorem, whose computer-assisted proof was accepted as a proof but never praised as beautiful. The connection between inevitability and explanation is close but not an identity: an explanatory proof whose route is entirely expected explains without surprising, and is by Lemma~\ref{l:inev} not beautiful.

\section{Rival accounts}

\emph{Economy and simplicity.} Hardy's third quality and the common identification of beauty with simplicity are, on the present account, conditions under which inevitability can be grasped. A proof must be short enough for the observer to hold its decisive feature in view, which is why long proofs are rarely called beautiful. But economy alone is not beauty: the shortest verification of an expected identity is economical and dull. \citet{InglisAberdein2015} found that simplicity and aesthetic appeal come apart in mathematicians' judgments, as the account predicts.

\emph{Enlightenment.} Rota's view captures inevitability and omits surprise. An enlightening proof of something one expected may deepen understanding without being beautiful, and Lemma~\ref{l:inev} shows that the omission matters. Rota is right that the aesthetic vocabulary in mathematics is tied to understanding; he is wrong that the tie makes it a misdescription. On Principle~\ref{p:rd}, the pleasure of grasping a surprising necessity is an aesthetic response like any other, and the evidence of \citet{ZekiEtAl2014} is consistent with this.

\emph{Unity in variety.} Hutcheson's formula is close to ours. Variety, in a theorem that unites apparently unrelated cases, is a source of surprise; uniformity, the single principle that comprehends them, is the source of inevitability. What the formula omits is the order of the two in the observer's experience. A result in which uniformity is visible from the start does not surprise, and one in which the variety is never reduced does not compel.

\section{Objections}

\begin{objection}[Familiarity]
Surprise fades with familiarity. An expert who has known Euclid's proof for decades is not surprised by it, yet she continues to find it beautiful. So surprise is not necessary.
\end{objection}

\begin{reply}
Beauty is a disposition to produce the response in qualified observers in suitable conditions, and Definition~\ref{d:obs} includes among suitable conditions the absence of prior familiarity with the route. The expert's judgment is grounded in her knowledge that the proof is disposed to surprise those who meet it properly, and in her capacity to reconstruct the surprise by recalling what the statement alone would lead one to expect. The same holds for a familiar piece of music whose modulations the listener knows. AES-02-a treats correctability in the same way: a judgment of beauty answers to the response of qualified observers, not to the judge's present response.
\end{reply}

\begin{objection}[Beautiful objects and definitions]
Mathematicians call objects, such as the icosahedron or the exceptional Lie group $E_8$, and definitions beautiful. These are not proofs and have no conclusions to be surprising or inevitable.
\end{objection}

\begin{reply}
Such ascriptions are grounded in what the object or definition makes inevitable. The icosahedron is admired for its symmetry, and what strikes the qualified observer is that a figure so constrained exists at all and that its properties follow from its symmetry group. A definition is called beautiful when the theorems that flow from it are surprising and follow from it with inevitability. The beauty of objects and definitions is derivative, as the beauty of theorems is, and Theorem~\ref{t:main} accounts for it in the same way.
\end{reply}

\begin{objection}[Disagreement]
Mathematicians disagree about which proofs are beautiful. If beauty were surprise plus inevitability, disagreement would reduce to disagreement about expectations, which should be easy to settle.
\end{objection}

\begin{reply}
Disagreements about beauty in mathematics are often disagreements about who counts as a qualified observer for a proof, since what is surprising depends on the background, and observers differ in background. A proof that surprises a number theorist may be routine to an analyst who knows the technique. The account predicts this source of disagreement and its partial resolution: when the background is specified, judgments converge. That some residual disagreement remains is consistent with AES-02-a, on which beauty is located in the relation between object and observer, and the observer kind may be specified more or less finely.
\end{reply}

\section{Conclusion}

Mathematical beauty is the combination of surprise and inevitability: the disposition of a proof to lead a qualified observer to a conclusion she could not have anticipated, by a route that shows it could not have been otherwise. Surprise without inevitability is curiosity, and inevitability without surprise is routine. Theorems, objects and definitions are beautiful through the proofs they make possible, and proofs that do not explain, among them most computer-assisted ones, are not beautiful.

\begin{thebibliography}{99}\small
\bibitem[Breitenbach(2015)]{Breitenbach2015} Breitenbach, A. (2015). Beauty in proofs: Kant on aesthetics in mathematics. \emph{European Journal of Philosophy}, 23(4), 955--977.
\bibitem[Hardy(1940)]{Hardy1940} Hardy, G.~H. (1940). \emph{A Mathematician's Apology}. Cambridge University Press.
\bibitem[Hutcheson(1725)]{Hutcheson1725} Hutcheson, F. (1725). \emph{An Inquiry into the Original of Our Ideas of Beauty and Virtue}. London.
\bibitem[Inglis and Aberdein(2015)]{InglisAberdein2015} Inglis, M. and Aberdein, A. (2015). Beauty is not simple: A multidimensional characterisation of mathematical aesthetics. \emph{Philosophia Mathematica}, 23(1), 87--109.
\bibitem[Montano(2014)]{Montano2014} Montano, U. (2014). \emph{Explaining Beauty in Mathematics: An Aesthetic Theory of Mathematics}. Springer.
\bibitem[Poincar\'e(1908)]{Poincare1908} Poincar\'e, H. (1908). \emph{Science et m\'ethode}. Flammarion.
\bibitem[Rota(1997)]{Rota1997} Rota, G.-C. (1997). The phenomenology of mathematical beauty. \emph{Synthese}, 111(2), 171--182.
\bibitem[Zeki et~al.(2014)]{ZekiEtAl2014} Zeki, S., Romaya, J.~P., Benincasa, D.~M.~T. and Atiyah, M.~F. (2014). The experience of mathematical beauty and its neural correlates. \emph{Frontiers in Human Neuroscience}, 8, 68.
\end{thebibliography}

\end{document}
